\chapter[Cartas dimensionales y transductores físicos]{Cartas angulares,
rectas dimensionales y transductores físicos comunes}
\label{chap:constantes-fisicas-genealogia-ampliacion-20260818}

Las constantes físicas de este capítulo no se introducen como parámetros
independientes. Todas reciben primero un objeto, una recta dimensional y un
transductor. La publicación en una carta metrológica se realiza al final. El
orden de dependencias es
\begin{equation}
 \label{eq:amp18-constantes-grafo-inicial}
 \begin{aligned}
 \mathrm{APP}\to\mathrm{TRIT}\to\mathrm{TPK}
 &\to(\pi_{\rm HMT},e_{\rm HMT},\varphi_{\rm HMT})
 \to\alpha_{\rm HMT},\\
 (\varphi_{\rm HMT},\alpha_{\rm HMT})
 &\to\Sigma_H\to10^{-34}\mathcal U_S
 \to\hbar_{\rm int},\\
 (\alpha_{\rm HMT},A,C^*,\Delta_4,q_\pm)
 &\to (V_{\rm CKM},Z_0,\mu_0,\varepsilon_0),\\
 (\hbar_{\rm int},t_0,\log3)
 &\to k_B\Theta_{\rm clk},\\
 (\hbar_{\rm int},t_0,c_{\rm int},\theta_G)
 &\to G_{\rm int}(\theta_G).
 \end{aligned}
\end{equation}
El Modelo Estándar, el Sistema Internacional de Unidades y las tablas
metrológicas se utilizan como lenguajes de publicación y contraste; no
determinan los estados ni los coeficientes de la construcción anterior.

\section{Los cuatro estratos de toda constante}

Sea \(\Lambda_{\rm D}\cong\mathbb Z^r\) la retícula libre de dimensiones y
\(\mathcal G_{\rm D}=\operatorname{Hom}(\Lambda_{\rm D},\mathbb R_{>0})\)
su grupo de cambios de base. Toda constante de esta Parte se expone en el
mismo orden:
\begin{enumerate}
 \item \textbf{constructor discreto o genealógico:} selecciona el objeto y
 su ruta sin cifra objetivo ni unidad;
 \item \textbf{caracterización analítica:} prueba existencia, unicidad,
 reversibilidad y sensibilidad dentro del dominio construido;
 \item \textbf{transducción dimensional:} lleva el invariante a una recta
 \(\mathcal L_d\) mediante el marco dimensional común;
 \item \textbf{comparación metrológica:} expresa la sección en una carta
 como SI, kelvin, newton, \(\mathrm{J\,s}\) o MeV.
\end{enumerate}

Un \emph{marco dimensional positivo} es una familia multiplicativa
\(\mathfrak u=(u_d)_{d\in\Lambda_{\rm D}}\), con
\(u_0=1\) y
\(u_{d+d'}=u_d\otimes u_{d'}\). El conjunto
\(\operatorname{Fr}_{\rm D}\) de estos marcos es un torsor de
\(\mathcal G_{\rm D}\) mediante
\begin{equation}
 \label{eq:amp18-constantes-torsor-marcos}
  (g\cdot\mathfrak u)_d:=g(d)u_d.
\end{equation}
En efecto, dados dos marcos \(\mathfrak u,\mathfrak v\), existe un único
carácter
\begin{equation}
 \label{eq:amp18-constantes-caracter-entre-marcos}
 g_{\mathfrak v/\mathfrak u}(d):=\frac{v_d}{u_d},
 \qquad
 g_{\mathfrak v/\mathfrak u}\cdot\mathfrak u=\mathfrak v.
\end{equation}
La multiplicatividad de ambos marcos implica
\(g(d+d')=g(d)g(d')\); por ello la acción es libre y transitiva. No se
postula un origen preferido en el torsor. Para cada dimensión
\(d\), los datos de realización comprenden un dominio
analítico \(X_d\) y un transductor equivariante
\begin{equation}
 \label{eq:amp18-constantes-transductor-equivariante}
 \mathfrak T_d:X_d\times\operatorname{Fr}_{\rm D}\longrightarrow
 \mathcal L_d,
 \qquad
 \mathfrak T_d(x,g\cdot\mathfrak u)
 =g(d)\,\mathfrak T_d(x,\mathfrak u).
\end{equation}
Si la caracterización analítica produce la coordenada adimensional
\(a_d(x)\), el transductor común es
\(\mathfrak T_d(x,\mathfrak u)=a_d(x)u_d\). Una realización queda
completamente tipada cuando determina \(X_d\), \(a_d\) y la
especialización de
\eqref{eq:amp18-constantes-transductor-equivariante}; el nombre de una
unidad, por sí solo, no define ese mapa.

\begin{theorem}[Separación entre invariante y coordenada metrológica]
\label{thm:amp18-constantes-cuatro-estratos}
Sea \(d\in\Lambda_{\rm D}\setminus\{0\}\). Ningún dato puramente
adimensional determina por sí solo una coordenada absoluta en
\(\mathcal L_d\) de manera natural respecto de todos
los cambios de unidad. Dados, además, el transductor equivariante
\(\mathfrak T_d\) de
\eqref{eq:amp18-constantes-transductor-equivariante} y un único marco
\(\mathfrak u\in\operatorname{Fr}_{\rm D}\), el constructor y la
caracterización determinan la sección
\(\mathfrak T_d(x,\mathfrak u)\in\mathcal L_d\). El teorema de
covariancia no crea ese transductor: demuestra su ley de cambio de marco.
Una carta metrológica publica después la coordenada de la sección; al
reexpresar la misma sección en otro marco, cambia el numeral y no la
magnitud representada.
\end{theorem}

\begin{proof}
Si \(F:X\to\mathcal L_d\) recibiera sólo datos adimensionales, el grupo
\(\mathcal G_{\rm D}\) actuaría trivialmente en \(X\). Para
\(g\in\mathcal G_{\rm D}=\operatorname{Hom}(\Lambda_{\rm D},
\mathbb R_{>0})\), denotemos por \(g(d)\) el valor de su carácter en
\(d\). La naturalidad implicaría \(F(x)=g(d)F(x)\) para todo
\(g\in\mathcal G_{\rm D}\). Como \(d\neq0\), existe \(g\) con
\(g(d)\neq1\), de donde \(F(x)=0\).
Esto prueba que el dato adimensional aislado no posee una sección natural
no nula. La existencia de \(\mathfrak T_d\) es, por tanto, un dato tipado necesario
de la realización, no una consecuencia del argumento de no existencia. El
marco \(\mathfrak u\) aporta una base positiva en cada generador dimensional
y \eqref{eq:amp18-constantes-transductor-equivariante} da la covariancia de
la salida. Para reexpresar una sección ya fijada bajo
\(u_d\mapsto g(d)u_d\), su coordenada cambia por
\(x\mapsto g(d)^{-1}x\), pero
\(xu_d=(g(d)^{-1}x)(g(d)u_d)\); la magnitud es la misma.
\end{proof}

Este teorema no reduce la derivación a una elección convencional. Obliga a
probar primero el invariante y a utilizar después una sola carta global. Una
unidad diferente para cada especie rompería la naturalidad y no pertenece a
la construcción.

\section{Datos comunes de la transducción}

El ángulo directo se obtiene de
\(A_{\rm deg}=1000\alpha_{\rm HMT}\). El contraángulo \(C^*\) se genera
después mediante la microfibra regional, el selector \(169\), \(D_A\),
\(H_5\) y la recurrencia determinantal. Una vez construido, la coordenada
angular elemental y la coordenada de retorno se definen localmente por
\begin{equation}
 \label{eq:amp18-constantes-alpha-angular-eta-retorno}
 \alpha_{\angle,\rm deg}:=\alpha_{\rm HMT}\,{}^\circ,
 \qquad
 \eta_{\rm ret}^{(\rm deg)}
 :=\frac{C^*_{\rm deg}}2-\alpha_{\angle,\rm deg}.
\end{equation}
Por tanto, y sólo después de la generación regional de \(C^*\), se obtiene
la identidad inversa
\begin{equation}
 \label{eq:amp18-constantes-contraangulo-inverso}
 \boxed{
 C^*_{\rm deg}
 =2\bigl(\eta_{\rm ret}^{(\rm deg)}+\alpha_{\angle,\rm deg}\bigr).}
\end{equation}
Esta igualdad reconoce y descompone el contraángulo; no es el selector que
lo genera.

El ángulo y el contraángulo se expresan en dos cartas coordinadas:
\[
 A_{\rm rad}=\frac{\pi A_{\rm deg}}{180},
 \qquad
 C^*_{\rm rad}=\frac{\pi C^*_{\rm deg}}{180}.
\]
Los canales orientados son
\begin{equation}
 \label{eq:amp18-constantes-qpm}
 q_-:=e^{-(A_{\rm rad}-C^*_{\rm rad})},
 \qquad
 q_+:=e^{-(A_{\rm rad}+C^*_{\rm rad})}.
\end{equation}
El mapa angular es reversible:
\begin{equation}
 \label{eq:amp18-constantes-qpm-inversa}
 A_{\rm rad}=-\frac12\log(q_+q_-),\qquad
 C^*_{\rm rad}=\frac12\log\!\left(\frac{q_-}{q_+}\right).
\end{equation}
La involución \(C^*\mapsto-C^*\) intercambia \(q_-\) y \(q_+\). El
invariante global que reaparece en los sectores de masa y mezcla es
\begin{equation}
 \label{eq:amp18-constantes-delta4}
 \Delta_4=\frac{\pi_{\rm HMT}-e_{\rm HMT}\log\pi_{\rm HMT}}{270}
 \left(1+\frac{\pi_{\rm HMT}}{729}\right).
\end{equation}
La sección de acción \(\hbar_{\rm int}\) procede del lector
\(\Sigma_H=\varphi_{\rm HMT}\alpha_{\rm HMT}^{16}\) y de su orden decimal
\(-34\); las unidades se asignan después mediante una base positiva de la
recta de acción.

\chapter[Constante de Boltzmann y transducción térmica]{Constante de
Boltzmann: transducción exacta entre información triádica, acción y energía}

\paragraph{Cuatro estratos de la constante de Boltzmann.}
\begin{enumerate}
 \item El constructor genealógico es la decisión TRIT de información
 \(\log3\), transportada por el reloj de orden \(108\).
 \item La caracterización analítica es la identidad
 \(k_B\Theta_{\rm clk}\log3
 =2\pi\hbar_{\rm int}/(108t_0)\) y su covariancia de escala.
 \item La transducción común es el morfismo positivo
 \(k_B:\mathcal L_\Theta\to\mathcal L_E\).
 \item La comparación metrológica comienza sólo al elegir el kelvin y el
 julio; la coordenada SI no selecciona el constructor.
\end{enumerate}

Sea \(\mathcal L_\Theta\) la recta positiva de temperatura y
\(\mathcal L_E\) la recta positiva de energía. La constante de Boltzmann es
el isomorfismo positivo
\begin{equation}
 \label{eq:amp18-constantes-boltzmann-morfismo}
 k_B:\mathcal L_\Theta\longrightarrow\mathcal L_E,
 \qquad [k_B]=\mathbf E-\boldsymbol\Theta.
\end{equation}
Esta definición separa temperatura y duración: la primera pertenece a
\(\mathcal L_\Theta\); la segunda, a la recta temporal
\(\mathcal L_T\).

Una decisión TRIT uniforme posee tres resultados visibles equiprobables y su
información es
\begin{equation}
 \label{eq:amp18-constantes-info-trit}
 I_{\rm TRIT}=\log3.
\end{equation}
Fijada una duración elemental \(t_0>0\), la periodicidad de orden \(108\)
define
\begin{equation}
 \label{eq:amp18-constantes-reloj-boltzmann}
 \omega_0=\frac{2\pi}{108t_0},
 \qquad
 \varepsilon_{\rm clk}=\hbar_{\rm int}\omega_0.
\end{equation}

\begin{definition}[Escala térmica de una decisión triádica]
La sección \(\Theta_{\rm clk}\in\mathcal L_\Theta^+\) es la única sección
que satisface
\begin{equation}
 \label{eq:amp18-constantes-boltzmann-estructural}
 \boxed{
 k_B\Theta_{\rm clk}\log3
 =\hbar_{\rm int}\frac{2\pi}{108t_0}.}
\end{equation}
Equivalentemente,
\[
 k_B\Theta_{\rm clk}
 =\frac{2\pi\hbar_{\rm int}}{108t_0\log3}.
\]
\end{definition}

La identidad \eqref{eq:amp18-constantes-boltzmann-estructural} determina la
energía por unidad de información triádica. Determina el producto
\(k_B\Theta_{\rm clk}\); la separación numérica de ambos factores requiere
elegir bases en las dos rectas.

\begin{proposition}[Covariancia de la escala térmica]
\label{prop:amp18-constantes-covariancia-termica}
Para todo \(\lambda>0\), la transformación
\[
 \Theta_{\rm clk}\longmapsto\lambda\Theta_{\rm clk},\qquad
 k_B\longmapsto\lambda^{-1}k_B
\]
conserva exactamente
\(k_B\Theta_{\rm clk}\log3\). Por consiguiente, antes de fijar una base de
\(\mathcal L_\Theta\), el invariante térmico es el morfismo
acción--información y su producto; el kelvin selecciona una coordenada de
ese morfismo, no el estado genealógico que lo origina.
\end{proposition}

\begin{proof}
El producto transformado es
\((\lambda^{-1}k_B)(\lambda\Theta_{\rm clk})=
k_B\Theta_{\rm clk}\). El miembro derecho de
\eqref{eq:amp18-constantes-boltzmann-estructural} no cambia, porque depende
de la sección de acción y del reloj ya construidos.
\end{proof}

En la carta SI, con \(1\,\mathrm J\) y \(1\,\mathrm K\) como bases, el
isomorfismo \eqref{eq:amp18-constantes-boltzmann-morfismo} tiene la coordenada
exacta
\begin{equation}
 \label{eq:amp18-constantes-kb-si}
 \boxed{
 k_B^{\rm SI}=1.380649\times10^{-23}\
 \mathrm{J\,K^{-1}}.}
\end{equation}
Esta cifra es definitoria en el SI vigente. La contribución estructural
consiste en construir el morfismo acción--información y el producto térmico
de \eqref{eq:amp18-constantes-boltzmann-estructural}.

Para una distribución no uniforme
\(p=(p_1,p_2,p_3)\), la cantidad \(\log3\) se reemplaza por
\[
 H(p)=-\sum_{j=1}^{3}p_j\log p_j.
\]
Si una evolución de rutas es biyectiva y conserva la medida, entonces
\(H(p)\) permanece invariante. En consecuencia, la contribución informacional
de Landauer satisface \(\Delta S=0\) y su cota inferior
\(Q\geq k_B\Theta\Delta S\) es nula. Esta afirmación se refiere al coste
mínimo asociado al borrado de información; no identifica ese límite con el
balance energético completo de cualquier realización dinámica.

\chapter[Impedancia, permeabilidad y permitividad del vacío]{Respuesta
constitutiva del vacío: impedancia, permeabilidad, permitividad y velocidad de
propagación}

\paragraph{Cuatro estratos de la respuesta constitutiva.}
\begin{enumerate}
 \item El constructor genealógico son los dos canales orientados
 \(q_\pm\), de los que se obtienen \(\mathcal E\) y \(\mathcal B\).
 \item La caracterización analítica son las identidades internas
 \(\widehat Z=\sqrt{\widehat\mu/\widehat\varepsilon}\) y
 \(\widehat c=(\widehat\mu\widehat\varepsilon)^{-1/2}\).
 \item La transducción común utiliza una única pareja de bases
 \((Z_{\rm ref},c_{\rm ref})\) y determina conjuntamente
 \(Z,\mu,\varepsilon,c\).
 \item La comparación metrológica publica después
 \(Z_0^{\rm HMT},\mu_0^{\rm HMT},\varepsilon_0^{\rm HMT}\) en la carta
 SI; ninguna de esas cifras interviene en \(q_\pm\).
\end{enumerate}

Los dos canales de \eqref{eq:amp18-constantes-qpm} producen un carácter par
\(\mathcal E\) y otro impar \(\mathcal B\):
\begin{equation}
 \label{eq:amp18-constantes-e-vacio}
 \mathcal E=
 \log[(1-q_+^{90})(1-q_-^{90})]
 -\log[(1-q_+^{120})(1-q_-^{120})],
\end{equation}
\begin{equation}
 \label{eq:amp18-constantes-b-vacio}
 \mathcal B=
 \log\frac{1-q_-^{90}}{1-q_+^{90}}
 -\log\frac{1-q_-^{120}}{1-q_+^{120}}.
\end{equation}
El intercambio de hojas conserva \(\mathcal E\) y cambia el signo de
\(\mathcal B\). Se definen las respuestas adimensionales
\begin{equation}
 \label{eq:amp18-constantes-respuestas-vacio}
 \widehat Z=e^{\mathcal B},\qquad
 \widehat c=e^{-\mathcal E},\qquad
 \widehat\mu=e^{\mathcal B+\mathcal E},\qquad
 \widehat\varepsilon=e^{-\mathcal B+\mathcal E}.
\end{equation}

\begin{theorem}[Identidades constitutivas internas]
Las cuatro respuestas de
\eqref{eq:amp18-constantes-respuestas-vacio} satisfacen exactamente
\begin{equation}
 \label{eq:amp18-constantes-identidades-vacio}
 \boxed{
 \sqrt{\frac{\widehat\mu}{\widehat\varepsilon}}=\widehat Z,
 \qquad
 (\widehat\mu\widehat\varepsilon)^{-1/2}=\widehat c.}
\end{equation}
\end{theorem}

\begin{proof}
Las definiciones dan
\[
 \frac{\widehat\mu}{\widehat\varepsilon}=e^{2\mathcal B},
 \qquad
 \widehat\mu\widehat\varepsilon=e^{2\mathcal E}.
\]
Todas las cantidades son positivas; al tomar la raíz positiva se obtiene
\eqref{eq:amp18-constantes-identidades-vacio}.
\end{proof}

La evaluación interna es
\[
 \widehat Z=0.9997245085389372\ldots,
 \qquad
 \widehat c=1.0002763104861575\ldots,
\]
\[
 \widehat\mu=0.9994483504793269\ldots,
 \qquad
 \widehat\varepsilon=0.9999992570966367\ldots.
\]
Las cuatro coordenadas contienen dos grados de libertad: el producto gobierna
la propagación común y el cociente conserva la orientación diferencial.

Sean \(Z_{\rm ref}>0\) y \(c_{\rm ref}>0\) bases de las rectas de
impedancia y velocidad. El transductor
\[
 Z=Z_{\rm ref}\widehat Z,
 \qquad c=c_{\rm ref}\widehat c,
 \qquad \mu=\frac Zc,
 \qquad \varepsilon=\frac1{Zc}
\]
conserva
\[
 Z=\sqrt{\mu/\varepsilon},
 \qquad c=(\mu\varepsilon)^{-1/2}.
\]

\subsection{Publicación de \texorpdfstring{\(Z_0\), \(\mu_0\) y
\(\varepsilon_0\)}{Z0, mu0 y epsilon0}}

En la carta SI posterior a 2019, \(q_e\), \(h\) y \(c\) son coordenadas
definitorias exactas; \(q_e\) denota la carga elemental y no el número de
Euler. A partir de la salida
\[
 \alpha_{\rm HMT}
 =0.007297352569283800997285105472380663\ldots
\]
se obtiene, en orden causal,
\begin{equation}
 \label{eq:amp18-constantes-vacio-si}
 \boxed{
 Z_0^{\rm HMT}=\frac{2h\alpha_{\rm HMT}}{q_e^2},
 \qquad
 \mu_0^{\rm HMT}=\frac{Z_0^{\rm HMT}}c,
 \qquad
 \varepsilon_0^{\rm HMT}=\frac1{Z_0^{\rm HMT}c}.}
\end{equation}
La evaluación es
\begin{equation}
 \label{eq:amp18-constantes-vacio-valores}
 \begin{aligned}
 Z_0^{\rm HMT}
 &=376.730313667167610257\ldots\ \Omega,\\
 \mu_0^{\rm HMT}
 &=1.256637062121047789\ldots\times10^{-6}\
   \mathrm{H\,m^{-1}},\\
 \varepsilon_0^{\rm HMT}
 &=8.854187812793002329\ldots\times10^{-12}\
   \mathrm{F\,m^{-1}}.
 \end{aligned}
\end{equation}
Por sustitución directa,
\begin{equation}
 \label{eq:amp18-constantes-velocidad-vacio}
 \boxed{
 (\mu_0^{\rm HMT}\varepsilon_0^{\rm HMT})^{-1/2}
 =299792458\ \mathrm{m\,s^{-1}}.}
\end{equation}
Así, impedancia, permeabilidad, permitividad y velocidad no constituyen cuatro
ajustes independientes: son publicaciones correlacionadas de dos caracteres
internos y de una misma salida \(\alpha_{\rm HMT}\).

\chapter[Constante gravitatoria]{Constante gravitatoria: transductor
dimensional, simetría dodecafásica y lectura holonómica}

\paragraph{Cuatro estratos de la constante gravitatoria.}
\begin{enumerate}
 \item El constructor genealógico circular es el reloj dodecafásico de
 orden \(108\); su cierre acción--Compton--gravedad selecciona
 \(\theta_{108}^{\rm circ}=(54/\pi)^2\).
 \item La caracterización analítica prueba la unicidad de ese selector. La
 expresión angular--holonómica constituye una segunda construcción
 homogénea y permanece separada mientras no se fije
 \(\mathcal L_{\rm hol}\).
 \item La transducción común es
 \(G_{\rm int}(\theta)=
 \theta c_{\rm int}^{5}t_0^2/\hbar_{\rm int}\).
 \item La carta decimal
 \(6.67430\times10^{-11}\) es una publicación declarada y no una
 evaluación retrospectiva de \(G_{\rm hol}\).
\end{enumerate}

La dimensión de la constante gravitatoria es
\[
 [G]=L^3M^{-1}T^{-2}.
\]
Sean
\[
 c_{\rm int}\in\mathcal L_C^+,
 \qquad t_0\in\mathcal L_T^+,
 \qquad \ell_0=c_{\rm int}t_0,
 \qquad \hbar_{\rm int}\in\mathcal L_S^+.
\]
Para \(\theta_G>0\), el transductor gravitatorio es
\begin{equation}
 \label{eq:amp18-constantes-transductor-g}
 \boxed{
 G_{\rm int}(\theta_G)
 =\theta_G\frac{c_{\rm int}^{5}t_0^2}{\hbar_{\rm int}}
 =\theta_G\frac{c_{\rm int}^{3}\ell_0^2}{\hbar_{\rm int}}.}
\end{equation}
Las dos expresiones poseen la dimensión correcta y son idénticas porque
\(\ell_0=c_{\rm int}t_0\). El transductor determina el tipo de \(G\); el
carácter \(\theta_G\) debe ser producido por una realización estructural.

\subsection{Selector dodecafásico de la carta de acción}

La realización circular de orden \(108\) se define por
\begin{equation}
 \label{eq:amp18-constantes-periodicidad-108}
 T_{108}=108t_0,
 \qquad
 \omega_{108}=\frac{2\pi}{108t_0},
 \qquad
 R_{108}=\frac{c_{\rm int}}{\omega_{108}}
 =\frac{54}{\pi}\ell_0.
\end{equation}
Para cualquiera de las dos secciones de acción \(\hbar_a\), sea
\[
 E_{108,a}=\hbar_a\omega_{108},
 \qquad
 m_{108,a}=\frac{E_{108,a}}{c_{\rm int}^{2}}.
\]
La longitud de Compton reducida cumple exactamente
\[
 \frac{\hbar_a}{m_{108,a}c_{\rm int}}=R_{108}.
\]

\begin{theorem}[Selector dodecafásico de cierre acción--Compton--gravedad]
Con la convención de radio gravitatorio reducido
\(r_g=Gm/c_{\rm int}^2\), la condición
\[
 r_g=R_{108}
\]
selecciona de manera única
\begin{equation}
 \label{eq:amp18-constantes-selector-108}
 \boxed{
 \theta_{108}^{\rm circ}=\left(\frac{54}{\pi}\right)^2.}
\end{equation}
En esta realización,
\(R_{108}=\bar\lambda_{\rm C}=r_g=\ell_{\rm P}^{\rm circ}\).
\end{theorem}

\begin{proof}
Al sustituir \eqref{eq:amp18-constantes-transductor-g} y
\(m_{108,a}=\hbar_a\omega_{108}/c_{\rm int}^2\), se obtiene
\[
 r_g
 =\frac{G_a(\theta)m_{108,a}}{c_{\rm int}^2}
 =\theta\,c_{\rm int}t_0^2\omega_{108}
 =\theta\frac{\pi}{54}\ell_0.
\]
La igualdad con \(R_{108}=(54/\pi)\ell_0\) implica
\(\theta=(54/\pi)^2\). La positividad garantiza unicidad.
\end{proof}

\subsection{Vía angular y holonómica}

Una segunda construcción parte de la fracción angular orientada
\begin{equation}
 \label{eq:amp18-constantes-delta-theta-g}
 \delta_\theta=\frac{A_{\rm deg}-C^*_{\rm deg}}{360}.
\end{equation}
Sea \(\widetilde{\mathcal K}_{27}\) un levantamiento enriquecido del retorno
de \(27\) pasos y sea
\[
 \Lambda_{27}=\mathcal L_{\rm hol}(\widetilde{\mathcal K}_{27})
\]
la longitud producida por un funcional geométrico de ruta. Si
\(\mathcal L_{\rm hol}\) es invariante por cambio cíclico de punto base, conjugación,
reetiquetado y subdivisión, la razón
\[
 \theta_G^{\rm hol}
 =\left(\delta_\theta
 \frac{\Lambda_{27}}{c_{\rm int}t_0}\right)^2
\]
produce la plantilla homogénea
\begin{equation}
 \label{eq:amp18-constantes-g-hol}
 \boxed{
 G_{\rm hol}
 =\frac{c_{\rm int}^{3}}{\hbar_{\rm int}}
 \bigl(\delta_\theta\Lambda_{27}\bigr)^2.}
\end{equation}
La homogeneidad de \eqref{eq:amp18-constantes-g-hol} es exacta. Su
especialización numérica exige construir \(\mathcal L_{\rm hol}\) y fijar una
normalización única e invariante bajo refinamiento. No se asigna aquí una
normalización a \(\Lambda_{27}\), ni se deducen cifras decimales de
\eqref{eq:amp18-constantes-g-hol} sin esa construcción.

\subsection{Carta decimal declarada}

La publicación decimal disponible se mantiene separada de la vía holonómica.
Para la terna tipada \([n,k,d]_3\), que no representa un numeral ternario,
se define
\begin{equation}
 \label{eq:amp18-constantes-lector-g}
 \mathcal R_G([n,k,d]_3;t,e_G)
 =10^{-n}\left(k+\frac{t}{10^3}+\frac{e_G}{10^d}\right).
\end{equation}
Con
\[
 [n,k,d]_3=[11,6,5]_3,
 \qquad t=674,
 \qquad e_G=30,
\]
la evaluación racional es
\begin{equation}
 \label{eq:amp18-constantes-g-decimal}
 \boxed{
 \mathcal R_G
 =10^{-11}\left(6+\frac{674}{1000}+\frac{30}{10^5}\right)
 =6.67430\times10^{-11}.}
\end{equation}
Tras elegir
\(\mathcal U_G=\mathrm{m^3\,kg^{-1}\,s^{-2}}\), la salida es
\(\mathcal R_G\mathcal U_G\). La comparación metrológica posterior es
\[
 G=6.67430(15)\times10^{-11}\
 \mathrm{m^3\,kg^{-1}\,s^{-2}}.
\]
La igualdad aritmética de \eqref{eq:amp18-constantes-g-decimal} pertenece a
la carta declarada; no se presenta como evaluación de
\eqref{eq:amp18-constantes-g-hol} ni se prolonga con decimales adicionales.

\chapter[Matriz CKM y violación CP]{Matriz CKM: ángulos de mezcla, fase de
violación CP e invariante de Jarlskog}

\paragraph{Cuatro estratos de las constantes de mezcla.}
\begin{enumerate}
 \item El constructor discreto es el sexteto de incidencia
 \((4,5,7,6,8,3)\), que determina las tres filas de
 \(M_{\rm CKM}\) antes de evaluar ángulos.
 \item La caracterización analítica es la invertibilidad de
 \(M_{\rm CKM}\), la relación exacta entre los cuatro ángulos y la
 unitariedad de \(V_{\rm CKM}\).
 \item Como estas constantes son adimensionales, la tercera capa es la
 carta angular común grados--radianes, aplicada una sola vez antes de las
 funciones trigonométricas.
 \item La comparación con el ajuste global se realiza después de obtener
 \(J\neq0\) y no selecciona ningún coeficiente de la matriz.
\end{enumerate}

El sexteto de incidencia
\[
 (b,p,q,h_6,o_8,t)=(4,5,7,6,8,3)
\]
determina las tres reglas sectoriales
\[
 R_{12}=(2,-p,bq),
 \qquad
 R_{23}=\left(0,q,-\frac{p^2}{2}\right),
 \qquad
 R_{13}=\left(1,-pb,-\frac{o_8-h_6}{t}\right).
\]
Por sustitución se obtiene la transformación racional
\begin{equation}
 \label{eq:amp18-constantes-mckm}
 \boxed{
 M_{\rm CKM}=
 \begin{pmatrix}
 2&-5&28\\
 0&7&-25/2\\
 1&-20&-2/3
 \end{pmatrix}.}
\end{equation}
Su determinante y su inversa son
\begin{equation}
 \label{eq:amp18-constantes-mckm-inversa}
 \det M_{\rm CKM}=-\frac{3857}{6}\neq0.
\end{equation}
Además,
\begin{equation}
 M_{\rm CKM}^{-1}=
 \begin{pmatrix}
 1528/3857&3380/3857&801/3857\\
 75/3857&176/3857&-150/3857\\
 6/551&-30/551&-12/551
 \end{pmatrix}.
\end{equation}
Los coeficientes se fijan antes de la evaluación angular.

Sea
\[
 h_{\rm deg}=\frac{C^*_{\rm deg}}6,
 \qquad
 \Delta_{\rm deg}=\frac{180}{\pi}\Delta_4.
\]
Los ángulos de mezcla y la fase se definen por
\begin{equation}
 \label{eq:amp18-constantes-angulos-ckm}
 \begin{pmatrix}\theta_{12}\\\theta_{23}\\\theta_{13}\end{pmatrix}
 =M_{\rm CKM}
 \begin{pmatrix}A_{\rm deg}\\h_{\rm deg}\\\Delta_{\rm deg}\end{pmatrix},
 \qquad
 \delta_{{\rm CP},\rm deg}=9A_{\rm deg}.
\end{equation}
Equivalentemente,
\begin{equation}
 \label{eq:amp18-constantes-ckm-cuatro-ecuaciones}
 \begin{aligned}
 \theta_{12}&=2A_{\rm deg}-5h_{\rm deg}+28\Delta_{\rm deg},\\
 \theta_{23}&=7h_{\rm deg}-\frac{25}{2}\Delta_{\rm deg},\\
 \theta_{13}&=A_{\rm deg}-20h_{\rm deg}
 -\frac{2}{3}\Delta_{\rm deg},\\
 \delta_{{\rm CP},\rm deg}&=9A_{\rm deg}.
 \end{aligned}
\end{equation}
La invertibilidad de \(M_{\rm CKM}\) da la relación exacta
\begin{equation}
 \label{eq:amp18-constantes-relacion-cp-angulos}
 3857\,\delta_{{\rm CP},\rm deg}
 =9\bigl(1528\theta_{12}+3380\theta_{23}+801\theta_{13}\bigr).
\end{equation}

La evaluación interna es
\begin{equation}
 \label{eq:amp18-constantes-valores-angulares-ckm}
 \boxed{\begin{aligned}
 \theta_{12}&=13.003313589509^\circ,&
 \theta_{23}&=2.398056157332^\circ,\\
 \theta_{13}&=0.213977382244^\circ,&
 \delta_{\rm CP}&=65.676173123554^\circ.
 \end{aligned}}
\end{equation}
En las funciones trigonométricas se usa siempre la carta en radianes:
\[
 s_{ij}=\sin\!\left(\frac{\pi\theta_{ij}}{180}\right),
 \qquad
 c_{ij}=\cos\!\left(\frac{\pi\theta_{ij}}{180}\right),
 \qquad
 \delta=\frac{\pi\delta_{{\rm CP},\rm deg}}{180}.
\]
La parametrización estándar es
\begin{equation}
 \label{eq:amp18-constantes-vckm}
 V_{\rm CKM}=
 \begin{pmatrix}
 c_{12}c_{13}&s_{12}c_{13}&s_{13}e^{-i\delta}\\
 -s_{12}c_{23}-c_{12}s_{23}s_{13}e^{i\delta}&
 c_{12}c_{23}-s_{12}s_{23}s_{13}e^{i\delta}&s_{23}c_{13}\\
 s_{12}s_{23}-c_{12}c_{23}s_{13}e^{i\delta}&
 -c_{12}s_{23}-s_{12}c_{23}s_{13}e^{i\delta}&c_{23}c_{13}
 \end{pmatrix}.
\end{equation}
Como producto de rotaciones unitarias elementales,
\(V_{\rm CKM}^{\dagger}V_{\rm CKM}=I_3\). Los módulos son
\begin{equation}
 \label{eq:amp18-constantes-modulos-vckm}
 |V_{\rm CKM}|\simeq
 \begin{pmatrix}
 0.974350259&0.225005836&0.003734601\\
 0.224873110&0.973489279&0.041841465\\
 0.008582400&0.041121745&0.999117283
 \end{pmatrix}.
\end{equation}

El invariante independiente de las redefiniciones de fase es
\begin{equation}
 \label{eq:amp18-constantes-jarlskog}
 J=c_{12}c_{23}c_{13}^2s_{12}s_{23}s_{13}\sin\delta.
\end{equation}

\begin{theorem}[Violación CP en la matriz construida]
Para los parámetros de
\eqref{eq:amp18-constantes-valores-angulares-ckm},
\begin{equation}
 \label{eq:amp18-constantes-jarlskog-valor}
 \boxed{J=3.1189723326632974\times10^{-5}\neq0.}
\end{equation}
Por tanto, \(V_{\rm CKM}\) no puede hacerse real mediante redefiniciones de
fase y describe violación de carga--paridad en el sector de tres
generaciones.
\end{theorem}

\begin{proof}
Los tres senos, los tres cosenos y \(\sin\delta\) son no nulos para los
valores de \eqref{eq:amp18-constantes-valores-angulares-ckm}. Su producto
es el valor de \eqref{eq:amp18-constantes-jarlskog-valor}.
\end{proof}

La comparación posterior con el ajuste global se resume en la tabla
siguiente. Las incertidumbres son marginales y los parámetros están
correlacionados; por ello, los cinco residuos no constituyen cinco pruebas
estadísticas independientes.
\begin{center}
\begin{tabular}{@{}lrrr@{}}
\toprule
Parámetro & HMT--MD & ajuste global & residuo normalizado\\
\midrule
\(\theta_{12}\) (grados)
&13.00331359&\(13.01287\pm0.03999\)&\(-0.24\)\\
\(\theta_{23}\) (grados)
&2.39805616&\(2.40082^{+0.04645}_{-0.03957}\)&\(-0.07\)\\
\(\theta_{13}\) (grados)
&0.21397738&\(0.215605^{+0.005042}_{-0.004756}\)&\(-0.34\)\\
\(\delta_{\rm CP}\) (grados)
&65.67617312&\(66.1193\pm1.4324\)&\(-0.31\)\\
\(10^5J\)
&3.11897233&\(3.16^{+0.13}_{-0.11}\)&\(-0.37\)\\
\bottomrule
\end{tabular}
\end{center}
El contraste no
interviene en \eqref{eq:amp18-constantes-mckm},
\eqref{eq:amp18-constantes-angulos-ckm} ni en la prueba de unitariedad.
La conclusión \(J\neq0\) se refiere a CP. La simetría CPT del registro actúa
sobre otro objeto y no se identifica con la matriz de mezcla.

\section{Coordinación causal de las derivaciones}

\begin{theorem}[Genealogía coordinada de las constantes físicas]
Fijados el estado enriquecido HMT, las cartas angulares y las bases positivas
de acción, duración, velocidad, temperatura y masa, se obtienen:
\begin{enumerate}
 \item el isomorfismo de Boltzmann y la identidad
 \(k_B\Theta_{\rm clk}\log3
 =\hbar_{\rm int}2\pi/(108t_0)\);
 \item las identidades constitutivas internas y las publicaciones de
 \eqref{eq:amp18-constantes-vacio-si},
 \[
  Z_0^{\rm HMT},\qquad
  \mu_0^{\rm HMT},\qquad
  \varepsilon_0^{\rm HMT};
 \]
 \item el transductor gravitatorio
 \eqref{eq:amp18-constantes-transductor-g}, el selector circular
 \eqref{eq:amp18-constantes-selector-108}, la plantilla holonómica
 \eqref{eq:amp18-constantes-g-hol} y la carta decimal
 \eqref{eq:amp18-constantes-g-decimal}, conservados como construcciones
 distintas;
 \item la transformación racional CKM, los cuatro parámetros angulares, la
 matriz unitaria y el invariante no nulo de Jarlskog.
\end{enumerate}
Cada salida se obtiene después de su genealogía. La publicación metrológica
reconoce la coordenada final sin intervenir en la selección del estado.
\end{theorem}

